
\chapter{Mostow rigidity theorem} \label{Mostow:chapter}
We have defined in Chapter \ref{Teichmuller:chapter} the
Teichm\"uller space $\Teich(S_g)$ of a genus-$g$ closed orientable surface $S_g$ as the space of all the hyperbolic metrics on $S_g$, considered up to isometries isotopic to the identity; we have then proved that $\Teich(S_g)$ is homeomorphic to $\matR^{6g-6}$ using the Fenchel--Nielsen coordinates.

This definition of $\Teich(M)$ actually applies to any closed hyperbolic manifold $M$, and we show here a striking difference between the dimensions two and three: if $\dim M = 3$ then $\Teich(M)$ is a single point. This strong result is known as the \emph{Mostow rigidity Theorem}.\index{Mostow--Prasad rigidity theorem} 

The impact of Mostow's rigidity on our understanding of three-dimensional topology cannot be overestimated. Thanks to this theorem every geometric information on a given closed hyperbolic three-manifold $M$ like its volume, geodesic spectrum, etc.~is promoted to a \emph{topological} invariant of $M$, that is it depends on the differentiable structure of $M$ only. In its strongest version, Mostow's rigidity says that the hyperbolic metric of $M$ is fully determined by the group $\pi_1(M)$ alone. 

We expose here Gromov's proof of Mostow's rigidity, which uses hyperbolic tetrahedra and introduces a nice invariant on closed manifolds of any dimension, called the \emph{simplicial volume}.\index{simplicial volume}

\section{Volume of tetrahedra} \label{volume:tetrahedra:section}
The volume of an ideal hyperbolic tetrahedron is a simple (but integral) formula that involves the \emph{Lobachevsky function}: we now prove this formula here. As a consequence we show that the regular ideal tetrahedron is the hyperbolic tetrahedron with maximal volume.

\subsection{The Lobachevsky function}
The \emph{Lobachevsky function} is\index{Lobachevsky function}
$$\Lambda(\theta) = -\int_0^\theta \log |2\sin t| dt.$$
The function $\log |2\sin t|$ is $-\infty$ on $\pi\matZ$ but is integrable, hence $\Lambda$ is well-defined and continuous on $\matR$. Its first derivatives are
$$\Lambda'(\theta) = -\log |2\sin \theta |, \qquad \Lambda''(\theta) = -\cot \theta.$$
The function $\Lambda$ has derivative $+\infty$ on $\pi\matZ$ and is an odd function, because its derivative is even. 

\begin{prop} \label{Loba:prop}
The function $\Lambda$ is $\pi$-periodic. We have $\Lambda (0) = \Lambda \left(\frac \pi 2 \right) = \Lambda (\pi) = 0$. The function $\Lambda$ is strictly positive on $\left(0,\frac \pi 2 \right)$, strictly negative on $\left(\frac \pi 2, \pi\right)$, and has absolute maximum and minimum at $\frac \pi 6 $ and $\frac 56 \pi$.
For all $m\in \matN$ the following holds:
$$\Lambda(m\theta) = m \sum_{k=0}^{m-1} \Lambda\left(\theta + \frac{k\pi}m\right).$$
\end{prop}
\begin{proof}
We prove the equality for $m=2$:
\begin{align*}
\frac{\Lambda(2\theta)}2 & = -\frac 12\int_0^{2\theta} \log |2\sin t|dt = - \int_0^\theta \log |2\sin 2t| dt \\
 & = - \int_0^\theta \log |2\sin t| dt - \int_0^\theta \log \left|2\sin \left(t+\frac \pi 2 \right) \right| dt \\
 & = \Lambda(\theta) - \int_{\frac \pi 2}^{\frac \pi 2 + \theta} \log |2 \sin t| dt \\
  & = \Lambda(\theta) + \Lambda\left(\theta + \frac \pi 2 \right) - \Lambda \left(\frac \pi 2 \right).
\end{align*}
By setting $\theta = \frac \pi 2$ we get $\Lambda(\pi)=0$. Since the derivative $\Lambda'$ is $\pi$-periodic and $\Lambda(\pi)=0$, also $\Lambda$ is $\pi$-periodic. Since $\Lambda$ is $\pi$-periodic and odd, we have $\Lambda \left(\frac \pi 2 \right) = 0$. We have also proved the formula for $m=2$.

To prove the formula for generic $m$ we use a generalisation of the duplication formula for the sinus. From the equality
$$z^m - 1 = \prod_{k=0}^{m-1} \left(z-e^{-\frac {2\pi ik}m}\right)$$
we deduce
$$2\sin(mt) = \prod_{k=0}^{m-1} 2\sin \left(t+\frac{k\pi}m\right)$$
and hence
\begin{align*}
\frac{\Lambda(m\theta)}m & = -\frac 1m \int_0^{m\theta} \log | 2\sin t| dt = - \int_0^\theta \log | 2 \sin (mt)| dt  \\
 & = - \sum_{k=0}^{m-1} \int_0^\theta \log \left|2 \sin \left( t + \frac {k \pi}m \right) \right| dt \\
 & = - \sum_{k=0}^{m-1} \left( \int_0^{\theta + \frac{k\pi}m} \log |2 \sin t| dt - \int_0^{\frac {k\pi}m} \log | 2\sin t | dt \right) \\
 & = - \sum_{k=0}^{m-1} \Lambda \left(\theta + \frac{k\pi}m \right) + C(m)
 \end{align*}
where $C(m)$ is a constant independent of $\theta$. By integrating both sides we get 
$$\frac 1m \int_0^\pi \Lambda(m\theta) = -\sum_{k=0}^{m-1}\int_0^\pi \Lambda\left(\theta + \frac{k\pi}m \right) + C(m)\pi.$$
Since $\Lambda$ is odd and $\pi$-periodic, we have
$$\int_0^\pi \Lambda (m\theta) = 0$$
for any integer $m$. Hence $C(m)=0$ and the formula is proved. Finally we note that $\Lambda''(\theta) = -\cot \theta $ is strictly negative in $(0,\frac \pi 2)$ and strictly positive in $(\frac \pi 2, \pi)$, hence $\Lambda$ is strictly positive in $(0, \frac \pi 2)$ and strictly negative in $(\frac \pi 2, \pi)$.
\end{proof}

\subsection{Volumes of ideal tetrahedra}
An ideal tetrahedron in $\matH^3$ is the convex hull of four non-planar ideal points. Quite surprisingly, every ideal tetrahedron has some non-trivial symmetries.\index{ideal tetrahedron} 

\begin{figure}
\begin{center}
\includegraphics[width = 4.5 cm] {\iftoggle{BW}{tetraedro4-BW}{tetraedro4}} 
\nota{Every pair of opposite edges in an ideal tetrahedron has an axis orthogonal to both which is a symmetry axis for the tetrahedron.}
\label{tetraedro4:fig}
\end{center}
\end{figure}

\begin{prop} \label{symmetries:ideal:tetrahedron:prop}
For any pair of opposite edges in an ideal tetrahedron $\Delta$ there is a unique line $r$ orthogonal to both as in Figure \ref{tetraedro4:fig} and $\Delta$ is symmetric with respect to a $\pi$-rotation around $r$.
\end{prop}
\begin{proof}
The opposite edges $e$ ed $e'$ are ultraparallel lines in $\matH^3$ and hence have a common perpendicular $r$. A $\pi$-rotation around $r$ inverts both $e$ and $e'$ but preserve the 4 ideal vertices of $\Delta$, hence $\Delta$ itself. 
\end{proof}


\begin{figure}
\begin{center}
\includegraphics[width = 10 cm] {\iftoggle{BW}{Loba-BW}{Loba}} 
\nota{The dihedral angles $\alpha, \beta, \gamma$ of an ideal tetrahedron. Opposite edges have the same angle and $\alpha+\beta+\gamma = \pi$ (left). To calculate the volume we use the half-space model, send a vertex to $\infty$, and divide the tetrahedron in six sub-tetrahedra (right).}
\label{Loba:fig}
\end{center}
\end{figure}

As a consequence, two opposite edges in $\Delta$ have coinciding dihedral angles as in Figure \ref{Loba:fig}-(left). Moreover, we have $\alpha + \beta + \gamma = \pi$ because a small horosphere based at a vertex intersects $\Delta$ into a Euclidean triangle with inner angles $\alpha, \beta$, and $\gamma$.
The regular ideal tetrahedron has of course equal angles $\alpha = \beta = \gamma = \frac \pi 3$.

\begin{teo} Let $\Delta$ be an ideal tetrahedron with dihedral angles $\alpha$, $\beta$ and $\gamma$. We have
$$\Vol(\Delta) = \Lambda(\alpha) + \Lambda(\beta) + \Lambda(\gamma).$$
\end{teo}
\begin{proof}
We represent $\Delta$ in the half-space model $H^3$ with one vertex $v_0$ at infinity and three vertices $v_1,v_2,v_3$ in $\matC$. Let $C$ be the circle containing $v_1$, $v_2$, and $v_3$: up to composing with elements in $\PSLC$ we can suppose that $C=S^1$. The Euclidean triangle $T\subset\matC$ with vertices $v_1$, $v_2$, and $v_3$ has interior angles $\alpha$, $\beta$, and $\gamma$. 

We first consider the case $0\in T$, that is $\alpha, \beta, \gamma \leqslant \frac \pi 2$. 
We decompose $T$ into six triangles as in Figure \ref{Loba:fig}: the tetrahedron $\Delta$ decomposes accordingly into six tetrahedra lying above them, and we prove that the one $\Delta_\alpha$ lying above the \iftoggle{BW}{light grey}{yellow} triangle has volume $\frac{\Lambda(\alpha)}2$. 
This proves the theorem.

The tetrahedron $\Delta_\alpha$ is the intersection of four half-spaces: three vertical ones bounded by the hyperplanes $y=0$, $x=\cos\alpha$, and $y=x\tan\alpha$, and one bounded by the half-sphere $z^2 = x^2+y^2$. Therefore
\begin{align*}
\Vol(\Delta_\alpha) & = \int_0^{\cos\alpha} dx \int_0^{x\tan \alpha} dy \int_{\sqrt{1-x^2-y^2}}^\infty \frac 1{z^3} dz \\
& = \int_0^{\cos\alpha} dx \int_0^{x\tan \alpha} dy \left[- \frac 1{2z^2} \right]_{\sqrt{1-x^2-y^2}}^\infty \\
& = \frac 12 \int_0^{\cos \alpha}dx \int_0^{x\tan\alpha}\frac 1{1-x^2-y^2}dy. 
\end{align*}
To solve this integral we use the relation 
$$\frac 1{1-x^2-y^2} = \frac 1{2\sqrt{1-x^2}} \left(\frac 1{\sqrt{1-x^2}-y} + \frac 1{\sqrt{1-x^2}+y}\right)$$
and hence $\Vol(\Delta_\alpha)$ equals
\begin{align*}
& \frac 14\int_0^{\cos\alpha} \!\!\!\!\!\! \frac {dx}{\sqrt{1-x^2}} 
\left(\left[-\log(\sqrt{1-x^2}-y)\right]_0^{x\tan \alpha} \!\!\!\!\! + \left[\log(\sqrt{1-x^2}+y)\right]_0^{x\tan\alpha}\right) \\
= & \frac 14 \int_0^{\cos\alpha} \!\!\!\!\! \frac{dx}{\sqrt{1-x^2}}
 \left(-\log(\sqrt{1-x^2}-x\tan\alpha)+\log(\sqrt{1-x^2}+x\tan\alpha)\right).
 \end{align*}
By writing $x=\cos t$ and hence $dx = -\sin t \, dt$ we obtain
\begin{align*} 
\Vol(\Delta_\alpha)  & = \frac 14 \int_{\frac\pi 2}^\alpha \frac{-\sin t}{\sin t} 
\left(-\log \frac{\sin t\cos\alpha - \cos t\sin \alpha}{\sin t \cos\alpha + \cos t \sin \alpha} \right) dt \\
 & = - \frac 14 \int_{\frac \pi 2}^\alpha \log \frac{\sin(t+\alpha)}{\sin(t-\alpha)}dt
 = - \frac 14 \int_{\frac \pi 2}^\alpha \log \frac{|2\sin(t+\alpha)|}{|2\sin(t-\alpha)|}dt \\
 & = \frac 14 \int_{2\alpha}^{\frac \pi 2+\alpha} \log | 2\sin t|dt - \frac 14 \int_0^{\frac \pi 2-\alpha} \log |2\sin t|dt \\
& = \frac 14\left(-\Lambda\left(\frac \pi 2 + \alpha \right) + \Lambda(2\alpha) + \Lambda\left(\frac \pi 2 - \alpha\right) \right) \\
& = \frac 14 \left(-\Lambda \left(\frac \pi 2 + \alpha \right) +2 \Lambda(\alpha) +2\Lambda\left(\frac \pi 2 + \alpha\right) - \Lambda\left(\frac \pi 2 + \alpha\right) \right) = \frac 12 \Lambda(\alpha)
\end{align*}
using Proposition \ref{Loba:prop}.

If $0\not \in T$ the triangle $T$ may be decomposed analogously into triangles, some of which contribute negatively to the volume, and we obtain the same formula.
\end{proof}

\begin{cor} \label{v3:cor}
The regular ideal tetrahedron is the hyperbolic tetrahedron of maximum volume.
\end{cor}
\begin{proof}
It is easy to prove that every hyperbolic tetrahedron is contained in an ideal tetrahedron: hence we may consider only ideal tetrahedra. Consider the triangle $T = \{0\leqslant \alpha, \beta, \alpha+\beta \leqslant \pi\}$ and 
\begin{align*}
f \colon \quad \ \ T & \longrightarrow  \matR \\ 
(\alpha, \beta) & \longmapsto \Lambda(\alpha) + \Lambda(\beta) + \Lambda(\pi-\alpha-\beta).
\end{align*}
The continuous function $f$ is null on $\partial T$ and strictly positive on the interior of $T$ because it measures the volume of the ideal tetrahedron of dihedral angles $\alpha, \beta, \gamma = \pi-\alpha-\beta$. Hence $f$ has at least a maximum on some interior point $(\alpha,\beta)$. The gradient 
$$
\nabla f = \begin{pmatrix}
\Lambda'(\alpha) - \Lambda'(\pi - \alpha-\beta) \\ \Lambda'(\beta) - \Lambda'(\pi-\alpha-\beta) 
\end{pmatrix}
 =
\begin{pmatrix}
-\log |2\sin \alpha| + \log |2\sin (\pi-\alpha-\beta)| \\ 
-\log |2\sin \beta| + \log |2\sin (\pi-\alpha-\beta) |
\end{pmatrix}
$$
must vanish there, and this holds if and only if $\sin\alpha = \sin (\pi-\alpha-\beta) = \sin \beta$, \emph{i.e.}~if and only if the tetrahedron has all dihedral angles equal to $\frac \pi 3$.
\end{proof}


\section{Simplicial volume}

Gromov has introduced a topological invariant on closed manifolds of any dimension called the \emph{simplicial volume}. This nice invariant can be used, among other things, to prove Mostow's rigidity theorem.\index{simplicial volume}

\subsection{Definition}
Gromov has introduced a measure of ``volume'' of a closed manifold $M$ which makes use only of the homology of $M$. Quite surprisingly, this notion of volume coincides (up to a factor) with the Riemannian one when $M$ is hyperbolic. 

Consider a topological space $X$ and its homology with ring $\matR$. We define the \emph{norm} of a cycle $\alpha = \lambda_1\alpha_1+\ldots +\lambda_h\alpha_h$ as follows:
$$|\alpha| = |\lambda_1|+\ldots +|\lambda_h|.$$

\begin{defn} The \emph{norm} of a class $a \in H_k(X,\matR)$ is the infimum of the norms of its elements: 
$$|a| = \inf \big\{|\alpha| \ \big|\ \alpha \in Z_k(X,\matR), [\alpha] = a \big\}.$$
\end{defn}

Recall that a \emph{seminorm} on a real vector space $V$ is a map $|\cdot|\colon V \to \matR_{\geqslant 0}$ such that
\begin{itemize}
\item $|\lambda v| = |\lambda||v|$ for any scalar $\lambda\in \matR$ and vector $v\in V$,
\item $|v+w| \leqslant |v|+|w|$ for any pair of vectors $v,w \in V$.
\end{itemize}
A \emph{norm} is a seminorm where $|v|=0$ implies $v=0$. The following is immediate.
\begin{prop} The norm $|\cdot |$ induces a seminorm on $H_k(X,\matR)$.
\end{prop}

Although it is only a seminorm, the function $|\cdot|$ is called a norm for simplicity. Let now $M$ be an oriented closed connected manifold: we know that $H_n(M,\matZ)\isom \matZ$ and the orientation of $M$ determines a \emph{fundamental class} $[M] \in H_n(M,\matZ)$ that generates the group. Moreover $H_n(M,\matR)\isom \matR$ and there is a natural inclusion\index{fundamental class}  
$$\matZ \isom H_n(M,\matZ) \hookrightarrow H_n(M,\matR) \isom \matR$$
hence the fundamental class $[M]$ is also an element of $H_n(M,\matR)$ and has a norm. 
\begin{defn} The \emph{simplicial volume} $\|M\|\in\matR_{\geqslant 0}$ of a closed oriented connected $M$ is the norm of its fundamental class:
$$\|M\| = |[M]|$$
\end{defn}
Since $|[M]| = |-[M]|$ the simplicial volume actually does not depend on the orientation. 
When $M$ is non-orientable we set $\|M\| = \big\|\tilde M\big\|/2$ where $\tilde M$ is the orientable double cover of $M$. The definition of $\|M\|$ is relatively simple but has various non-obvious consequences.

\subsection{Properties}
A continuous map $f\colon M \to N$ between closed oriented $n$-manifolds induces a homomorphism $f_*\colon H_n(M,\matZ) \to H_n(N,\matZ)$, and recall that the \emph{degree} of $f$ is the integer $\deg f$ such that
$$f_*([M]) = \deg f \cdot [N].$$

\begin{prop} \label{grado:prop}
Let $f\colon M\to N$ be a continuous map between closed oriented manifolds. The following inequality holds:
$$\|M\| \geqslant |\deg f| \cdot \|N\|.$$
\end{prop}
\begin{proof}
Every description of $[M]$ as a cycle $\lambda_1\alpha_1+\ldots +\lambda_h\alpha_h$ induces a description of $f_*([M]) = \deg f[N]$ as a cycle $\lambda_1f\circ\alpha_1+\ldots +\lambda_hf\circ\alpha_h$ with the same norm (or less, if there is some cancelation). 
\end{proof}

\begin{cor}
If $M$ and $N$ are closed orientable and homotopically equivalent $n$-manifolds then $\|M\|=\|N\|$.
\end{cor}
\begin{proof}
A homotopic equivalence consists of two maps $f\colon M\to N$ and $g\colon N \to M$ whose compositions are both homotopic to the identity. In particular both $f$ and $g$ have degree $\pm 1$.
\end{proof}

\begin{cor}
If $M$ admits a continuous self-map $f\colon M \to M$ of degree at least two then $\|M\| = 0$.
\end{cor}
\begin{cor}
Every sphere $S^n$ has norm zero. More generally we have $\|M\times S^n\|=0$ for every closed $M$ and any $n\geqslant 1$.
\end{cor}
\begin{proof}
A sphere $S^n$ admits self-maps of degree $\geqslant 2$, and hence also the product $M\times S^n$ does.
\end{proof}
Among the genus-$g$ surfaces $S_g$, we deduce that the sphere and the torus have simplicial volume zero. We will see soon that every surface of genus $g\geqslant 2$ has positive simplicial volume. When the continuous map is a covering the inequality from Proposition \ref{grado:prop} is promoted to an equality.

\begin{prop}
If $f\colon M \to N$ is a degree-$d$ covering we have 
$$\|M\| = d \cdot \|N \|.$$
\end{prop}
\begin{proof} 
The reason for this equality is that cycles can be lifted and projected through the covering. More precisely, we already know that $\|M\|\geqslant d\cdot \|N\|$. Conversely, let $\alpha = \lambda_1\alpha_1+\ldots +\lambda_h\alpha_h$ represent $[N]$; each $\alpha_i$ is a map $\Delta_n\to N$. Since $\Delta_n$ is simply connected, the map $\alpha_i$ lfts to $d$ distinct maps $\alpha_i^1,\ldots, \alpha_i^d\colon \Delta_n\to N$. The chain $\tilde\alpha = \sum_{ij} \lambda_i\alpha_i^j$ is a cycle in $M$ and $f_*(\tilde \alpha) = d\alpha$. Hence $\|M\|\leqslant d\cdot \|N\|$.
\end{proof}

We also note the following fact.

\begin{prop}
If $M$ is triangulated with $k$ simplices, then $\|M\| \leqslant k$.
\end{prop}
\begin{proof}
Up to taking a double cover we may suppose that $M$ is oriented. The closed $n$-manifold $M$ is triangulated into some simplices $\Delta_1, \ldots, \Delta_k$, and
we fix an orientation-preserving parametrisation $s_i\colon \Delta \to \Delta_i$ of each.
We would like to say that $s_1+\ldots + s_k$ is a fundamental cycle, however this singular chain is not necessarily a cycle because the restriction to adjacent faces coincide only up to the symmetries $S_{n+1}$ of $\Delta$. 

We can fix this problem easily by averaging each $s_i$ on all its permutations, that is we substitute each $s_i$ with $\frac{1}{(n+1)!}\sum_{\sigma \in S_{n+1}} (-1)^{{\rm sgn}(\sigma)} s_i \circ \sigma$. Now $s=s_1+\ldots + s_k$ is a fundamental cycle and $|s| = k$.
\end{proof}

\subsection{Seifert manifolds}
We can now calculate the simplicial volume of Seifert manifolds.

\begin{prop}
If $M$ is a closed Seifert 3-manifold then $\|M\|=0$.
\end{prop}
\begin{proof}
Every Seifert manifold is finitely covered by a product $S \times S^1$ (if $e=0$) or by a bundle $\big(S,(1,1)\big)$ with Euler number 1 (if $e\neq 0$). In the first case we are done since $\|S\times S^1\|=0$. 

In the second, if $S=S^2$ then $\big(S^2,(1,1)\big)=S^3$ and $\|S^3\|=0$, so we suppose $\chi(S)\leqslant 0$. There is a universal $K>0$ such that $\big(S,(1,1)\big)$ triangulates with at most $K|\chi (S)|+K$ tetrahedra (exercise) and therefore $\big\|\big(S,(1,1)\big)\big\| \leqslant K|\chi(S)|+K$. Exercise \ref{e:construct:ex} shows that for every $e>0$ there is a degree-$e$ covering $\tilde S \to S$ and two degree-$e$ coverings
$$\big(\tilde S, (1,1)\big) \longrightarrow \big(\tilde S, (1,e)\big) \longrightarrow \big(S,(1,1)\big)$$
which compose to a degree-$e^2$ covering 
$\big(\tilde S, (1,1)\big) \to \big(S,(1,1)\big)$. Therefore
$$\big\|\big(S, (1,1) \big) \big\| = \frac{\big\|\big(\tilde S, (1,1) \big) \big\|}{e^2} \leqslant \frac{K|\chi(\tilde S)| + K}{e^2} \leqslant \frac{Ke|\chi(S)| + K}{e^2} \to 0$$
as $e \to \infty$.
\end{proof}


\subsection{Simplicial and hyperbolic volume}
In the next pages we will prove the following theorem. Let $v_3$ be the volume of the regular ideal tetrahedron in $\matH^3$.

\begin{teo} \label{Gromov:teo}
Let $M$ be a closed hyperbolic $3$-manifold. We have 
$$\Vol(M) = v_3\|M\|.$$
\end{teo}
The theorem furnishes in particular some examples of manifolds with positive simplicial volume and shows that $\Vol(M)$ is a topological invariant of $M$, thus generalising the Gauss-Bonnet theorem to dimension $n=3$. Mostow rigidity will then strengthen this result in dimension $n = 3$, showing that the hyperbolic metric itself (not only its volume) is a topological invariant.

Both quantities $\Vol(M)$ and $\|M\|$ are multiplied by $d$ if we substitute $M$ with a degree-$d$ covering. In particular, up to substituting $M$ with its orientable double cover we can suppose that $M$ is orientable. 

\subsection{Cycle straightening}
The \emph{straight singular $k$-simplex} with vertices $v_0,\ldots, v_{k} \in \matH^n$ is the map
\begin{align*}
\alpha\colon \quad \ \Delta_k & \longrightarrow \matH^n \\
(t_0,\ldots, t_{k}) & \longmapsto t_0v_0 + \ldots + t_{k}v_{k}
\end{align*}
defined using convex combinations, see Section \ref{convex:subsection}. If the $k+1$ vertices are not contained in a $(k-1)$-plane the singular $k$-simplex is \emph{non-degenerate} and its image is a hyperbolic $k$-simplex.

The \emph{straightening} $\alpha^\st$ of a singular simplex $\alpha\colon \Delta_k \to \matH^n$ is the straight singular simplex with the same vertices of $\alpha$. The straightening $\alpha^\st$ of a singular simplex $\alpha\colon\Delta_k\to M$ in a hyperbolic manifold $M=\matH^n/_\Gamma$ is defined by lifting the singular simplex in $\matH^n$, straightening it, and projecting it back to $M$ by composition with the covering map.
Different lifts produce the same straightening in $M$ because they are related by isometries of $\matH^n$. 

The straightening extends by linearity to a homomorphism
$$\st\colon C_k(M,R) \to C_k(M,R)$$
which commutes with $\partial$ and hence induces a homomorphism in homology
$$\st_*\colon H_k(M,R) \to H_k(M,R).$$
\begin{prop} The map $\st_*$ is the identity.
\end{prop}
\begin{proof}
We may define a homotopy between a singular simplex $\sigma$ and its straightening $\sigma^{\rm st}$ using the convex combination
$$\sigma^t(x) =t\sigma(x) + (1-t)\sigma^{\rm st}(x).$$
This defines a chain homotopy between $\st_*$ and $\id$ via the same technique used to prove that homotopic maps induce the same maps in homology.
\end{proof}

The \emph{abstract volume} of a straightened singular $n$-simplex $\alpha\colon\Delta_n \to M$ is the volume of its lift in $\matH^n$ and may also be calculated as
$$\left|\int_\alpha \omega\right|$$
where $\omega$ is the volume form on $M$ pulled back along $\alpha$. 
If $\alpha$ is non-degenerate, we say that its \emph{sign} is \emph{positive} if $\alpha$ is orientation-preserving and \emph{negative} otherwise: equivalently, it is the sign of $\int_\alpha\omega$.

We now concentrate on the dimension $n=3$ where we know that a tetrahedron has maximum volume $v_3$ if and only if it is both regular and ideal: a straight simplex is compact and hence its abstract volume is strictly smaller than $v_3$. 
We can now easily prove one inequality. 

\begin{prop} \label{Gromov:1:prop}
Let $M$ be a closed hyperbolic $3$-manifold. We have
$$\Vol(M) \leqslant v_3 \|M\|.$$
\end{prop}
\begin{proof}
As we said above, we can suppose that $M$ is orientable. Take a cycle $\alpha = \lambda_1\alpha_1+\ldots+\lambda_k\alpha_k$ that represents $[M]$. We can suppose it is straightened, because the straightening preserves both the coefficients and the homology class. Let $\omega$ be the volume form on $M$. We get 
$$\Vol (M) = \int_M \omega = \int_\alpha \omega = \lambda_1 \int_{\alpha_1}\omega + \ldots + \lambda_k \int_{\alpha_k}\omega.$$
The quantity $\big|\int_{\alpha_i}\omega\big|$ is the abstract volume of $\alpha_i$. Hence $\big|\int_{\alpha_i}\omega\big|<v_3$ and
$$\Vol(M) < \big(|\lambda_1|+ \ldots +|\lambda_k|\big)v_3.$$
This holds for all $\alpha$, hence $\Vol(M) \leqslant v_3 \|M\|$.
\end{proof}

The proof of the converse inequality is less immediate.

\subsection{Efficient cycles} \label{efficienti:subsection}
Let $M=\matH^3/_\Gamma$ be a closed oriented hyperbolic 3-manifold. 
An \emph{$\varepsilon$-efficient cycle} for $M$ is a straightened cycle
$$\alpha = \lambda_1\alpha_1 + \ldots + \lambda_k\alpha_k$$
representing $[M]$ where the abstract volume of $\alpha_i$ if bigger than $v_3-\varepsilon$ and the sign of $\alpha_i$ is coherent with the sign of $\lambda_i$, for all $i$.

We will construct an $\varepsilon$-efficient cycle for every $\varepsilon >0$. This will conclude the proof of Theorem \ref{Gromov:teo} in virtue of the following:

\begin{lemma}
If for every $\varepsilon>0$ the manifold $M$ admits an $\varepsilon$-efficient cycle, then we have $\Vol(M) \geqslant v_3 \|M\|$.
\end{lemma}
\begin{proof}
Let $\alpha = \lambda_1\alpha_1 + \ldots \lambda_k\alpha_k$ be an $\varepsilon$-efficient cycle and $\omega$ be the volume form on $M$. Coherence of signs gives $\lambda_i\int_{\alpha_i}\omega>0$ for all $i$.
We get
\begin{align*}
\Vol(M) & = \int_M \omega = \int_\alpha \omega = \lambda_1 \int_{\alpha_1} \omega + \ldots + \lambda_k \int_{\alpha_k} \omega  \\
& \geqslant \big(|\lambda_1| + \ldots + |\lambda_k|\big) \cdot (v_3 - \varepsilon).
\end{align*}
Therefore $\Vol(M) \geqslant \|M\| \cdot (v_3 - \varepsilon)$ for all $\varepsilon > 0$.
\end{proof}

It remains to construct $\varepsilon$-efficient cycles.

\begin{ex} \label{volume:continuo:ex}
If $\Delta^i$ is a sequence of tetrahedra in $\matH^3$ whose vertices tend to the vertices of a regular ideal tetrahedron in $\partial\matH^3$, then
$$\Vol(\Delta^i)\to v_3.$$
\end{ex}

For any $t>0$, let $\Delta(t)$ be a regular tetrahedron obtained as in Section \ref{platonici:subsection} as follows.
Pick a point $x\in \matH^3$ and a regular tetrahedron in the Euclidean tangent space $T_x\matH^3$, centred at the origin with vertices at distance $t$ from it, project the vertices in $\matH^3$ via the exponential map, and pick their convex hull.

A \emph{$t$-simplex} is a tetrahedron isometric to $\Delta(t)$ equipped with an ordering of its vertices: the ordering allows us to consider it as a straightened singular simplex. Let $S(t)$ be the set of all $t$-simplices in $\matH^3$.

\begin{ex} The group $\Iso(\matH^3)$ acts on $S(t)$ freely and transitively.
\end{ex}

Recall from Corollary \ref{unimodular:cor} that $\Iso(\matH^3)$ is unimodular:
the Haar measure on $\Iso(\matH^3)$ induces an $\Iso(\matH^3)$-invariant measure on $S(t)$. 

Let $M=\matH^3/_\Gamma$ be a closed hyperbolic 3-manifold and $\pi\colon\matH^3 \to M$ the covering projection. Fix a base point $x_0 \in \matH^3$ and consider its orbit $O=\Gamma x_0$. Consider the set
$$\Sigma = \Gamma^{4}/_\Gamma$$
of the $4$-uples $(g_0,g_1,g_2, g_3)$ considered up to the diagonal action of $\Gamma$:
$$g \cdot (g_0,g_1,g_2, g_3) = (g g_0,gg_1, gg_2, g g_3).$$
An element $\sigma = (g_0,g_1,g_2, g_3) \in \Sigma$ determines a singular simplex $\tilde\sigma$ in $\matH^3$ with vertices $g_0(x_0), g_1(x_0), g_2(x_0), g_3(x_0)\in O$ only up to translations by $g\in\Gamma$, hence it gives a well-defined singular simplex in $M$, which we still denote by $\sigma$. We now introduce the chain
$$\alpha(t) = \sum_{\sigma \in \Sigma} \lambda_\sigma(t) \cdot \sigma$$
for some appropriate real coefficients $\lambda_\sigma(t)$ that we now define. The base point $x_0$ determines the Dirichlet tessellation of $\matH^3$ into domains $D(g(x_0))$, $g\in\Gamma$. For $\sigma = (g_0,g_1,g_2, g_3)$ we let $S^+_\sigma(t)\subset S(t)$ be the set of all positive $t$-simplices whose $i$-th vertex lies in $D(g_i(x_0))$ for all $i$. The number $\lambda^+_\sigma(t)$ is the measure of $S^+_\sigma(t)$. We define analogously $\lambda^-_\sigma(t)$ and set
$$\lambda_\sigma(t) = \lambda^+_\sigma(t) - \lambda^-_\sigma(t).$$

\begin{lemma} \label{ciclo:lemma}
The chain $\alpha(t)$ has finitely many addenda and is a cycle. If $t$ is sufficiently big the cycle $\alpha(t)$ represents a positive multiple of $[M]$ in the group $H_3(M,\matR)$.
\end{lemma}
\begin{proof}
We prove that the sum is finite. Let $d, T$ be the diameters of $D(x_0)$ and of a $t$-simplex. 
We write $\sigma = (\id,g_1,g_2,g_3)$ for all $\sigma \in \Sigma$: that is, all simplices have their first vertex at $x_0$. 
If $\lambda_\sigma(t)\neq 0$ then $d(g_ix_0, x_0) < 2d+T$ for all $i$: therefore $\alpha(t)$ has finitely many addenda (because $O$ is discrete).

We prove that $\alpha(t)$ is a cycle. The boundary $\partial \alpha(t)$ is a linear combination of straight $2$-simplices with vertices in $(g_0x_0, g_1x_0, g_2x_0)$ as $g_0, g_1$, and $g_2$ vary. The coefficient of one such $2$-simplex is
$$\sum_{g\in\Gamma} \big(- \lambda_{(g,g_0,g_1,g_2)}(t) + \lambda_{(g_0,g,g_1,g_2)}(t)
- \lambda_{(g_0,g_1,g,g_2)}(t) + \lambda_{(g_0,g_1,g_2,g)}(t)\big).$$
We prove that each addendum summed along $g\in \Gamma$ is zero; for simplicity we consider the last addendum and get
$$\sum_{g\in\Gamma} \lambda_{(g_0, g_1, g_{2}, g)}(t) = \sum_{g\in\Gamma} \lambda_{(g_0, g_1, g_{2}, g)}(t)^+ - \sum_{g\in\Gamma} \lambda_{(g_0, g_1, g_{2}, g)}(t)^-.$$
The first addendum measures the positive $t$-simplices whose first $3$ vertices lie in $D(g_0(x_0)), \ldots, D(g_2(x_0))$, the second measures the negative $t$-simplices with the same requirement. These two subsets have the same volume in $S(t)$ because they are related by the involution $r\colon S(t) \to S(t)$ that mirrors a simplex with respect to its first facet. 

We show that for sufficiently big $t$ the cycle is a positive multiple of $[M]$. Let $t$ be sufficiently big so that two vertices in a $t$-simplex have distance 
bigger than $2d$. This condition implies that if there is a positive $t$-simplex with vertices in $D(g_0(x_0)), \ldots, D(g_3(x_0))$, then any straight simplex with vertices in $D(g_0(x_0)), \ldots, D(g_3(x_0))$ is positive. 
Therefore in the expression
$$\alpha(t) = \sum_{\sigma \in \Sigma} \lambda_\sigma(t) \cdot \sigma$$
the signs of $\lambda_\sigma(t)$ and $\sigma$ are coherent and
$$\int_{\alpha(t)} \omega = \sum_{\sigma\in \Sigma} \lambda_\sigma(t) \cdot \int_{\sigma}\omega > 0.$$
Therefore $\alpha(t)$ is a positive multiple of $[M]$.
\end{proof}

For sufficiently big $t$ we have $\alpha(t) = k_t[M]$ in homology for some $k_t>0$. The rescaled $\bar\alpha(t) = \alpha(t)/k_t$ hence represents $[M]$.
We have found our $\varepsilon$-efficient cycles. 

\begin{lemma} \label{efficiente:lemma}
For any $\varepsilon>0$ there is a $t_0>0$ such that $\bar\alpha(t)$ is $\varepsilon$-efficient for all $t>t_0$.
\end{lemma}
\begin{proof}
Let $d$ be the diameter of the Dirichlet domain $D(x_0)$. Let a \emph{quasi $t$-simplex} be a simplex whose vertices are at distance $<d$ from those of a $t$-simplex. By construction $\bar\alpha(t)$ is a linear combination of quasi $t$-simplices.

We now show that for any $\varepsilon>0$ there is a $t_0>0$ such that for all $t>t_0$ every quasi $t$-simplex has volume bigger than $v_3-\varepsilon$. By contradiction, let $\Delta^t$ be a sequence of quasi $t$-simplices of volume smaller than $v_3-\varepsilon$ with $t\to\infty$. The vertices of $\Delta^t$ are $d$-closed to a $t$-simplex $\Delta_*^t$, and we move the pair $\Delta^t, \Delta^t_*$ isometrically so that the $t$-simplices $\Delta_*^t$ have the same barycenter. Now both the vertices of $\Delta^t$ and $\Delta^t_*$ tend to the vertices of an ideal regular tetrahedron and Exercise \ref{volume:continuo:ex} gives a contradiction.
\end{proof}

The previous lemmas together prove the second half of Theorem \ref{Gromov:teo}.
\begin{cor}
Let $M$ be a closed hyperbolic $3$-manifold. We have
$$\Vol(M) \geqslant v_3 \|M\|.$$
\end{cor}

Theorem \ref{Gromov:teo} has some non-trivial consequences.

\begin{cor} \label{degree:d:cor}
Let  $M, N$ be two closed orientable hyperbolic $3$-manifolds. If there is a map $f\colon M\to N$ of degree $d$ then $\Vol(M)\geqslant |d| \cdot \Vol(N)$.
\end{cor}

\begin{cor} \label{stesso:volume:cor}
Two homotopically equivalent and closed hyperbolic 3-manifolds have the same volume.
\end{cor}

\begin{rem}
If we are able to prove that the regular ideal simplex has the maximum volume $v_n$ among hyperbolic $n$-simplices in $\matH^n$, then the whole proof extends as is from $\matH^3$ to $\matH^n$ and shows that $\Vol(M) = v_n \cdot |M|$ for any closed hyperbolic $n$-manifold $M$. This is obviously true when $n=2$, and we get $v_2 = \pi$. 
\end{rem}

\section{Mostow rigidity}

\subsection{Introduction}
We want to prove the following.
\begin{teo}[Mostow rigidity] \label{Mostow:teo}
Let $M$ and $N$ be two closed connected orientable hyperbolic 3-manifolds. Every isomorphism $\pi_1(M) \stackrel \sim\to \pi_1(N)$ between fundamental groups is induced by a unique isometry $M\stackrel \sim\to N$.
\end{teo}
This very powerful theorem says that an algebraic isomorphism between fundamental groups alone is enough to produce and characterise an isometry. 

\begin{cor} Two closed orientable hyperbolic 3-manifolds with isomorphic fundamental groups are isometric.
\end{cor}

We prove Mostow's rigidity in this section. All the ingredients are already there, we only need to make a little last effort.

We note that closed hyperbolic manifolds are aspherical because their universal cover $\matH^n$ is contractible. For such manifolds every isomorphism $\pi_1(M)\to \pi_1(N)$ is induced by a homotopy equivalence $f\colon M \to N$, unique up to homotopy: see Corollary \ref{aspherical:cor}. To prove Mostow's theorem we need to promote this homotopy equivalence to an isometry in dimension $n=3$. We already know that $\Vol(M) = \Vol(N)$ by Corollary \ref{stesso:volume:cor}.

\subsection{Proof of Mostow's theorem}
Let $f\colon M \to N$ be a smooth homotopy equivalence. Recall from Theorem \ref{estensione:teo} that $f$lifts to a map $\tilde f\colon \matH^3 \to \matH^3$ which extends continuously to a homeomorphism $\tilde f\colon\partial \matH^3 \to \partial\matH^3$ of the boundary spheres.
We start with a lemma.

\begin{lemma} \label{regolare:lemma}
The extension $\tilde f\colon \partial\matH^3 \to \partial \matH^3$ sends the vertices of every regular ideal simplex to the vertices of some regular ideal simplex. 
\end{lemma}
\begin{proof}
Let $w_0,\ldots, w_3$ be vertices of a regular ideal simplex and suppose by contradiction that their images $\tilde f(w_0),\ldots ,\tilde f(w_3)$ span a non-regular ideal simplex, which has volume smaller than $v_3-2\delta$ for some $\delta >0$. 
By Exercise \ref{volume:continuo:ex} there are neighbourhoods $U_i$ of $w_i$ in $\overline{\matH^3}$ for $i=0,\ldots, 3$ such that the volume of the simplex with vertices $\tilde f(u_0), \ldots, \tilde f(u_3)$ is smaller than $v_3-\delta$ for any choice of $u_i\in U_i$. 

In Section \ref{efficienti:subsection} we have defined a cycle
$$\alpha(t) = \sum_{\sigma \in \Sigma} \lambda_\sigma(t) \cdot \sigma$$
where $t$ depends on $\varepsilon$. We say that a singular simplex $\sigma\in \Sigma$ is \emph{bad} if its $i$-th vertex is contained in $U_i$ for all $i$. 
Let $\Sigma^{\rm bad}\subset\Sigma$ be the subset of all bad singular simplices and define
$$\alpha(t)^{\rm bad} = \sum_{\sigma \in \Sigma^{\rm bad}} \lambda_\sigma(t) \cdot \sigma.$$
We want to estimate $|\alpha(t)|$ and $|\alpha(t)^{\rm bad}|$. We prove that  
$$|\alpha(t)| = \sum_{\sigma\in \Sigma} |\lambda_\sigma(t)|$$
is a real number independent of $t$: let $S_0\subset S(t)$ be the set of all $t$-simplices having the first vertex in the Dirichlet domain $D(x_0)$ of the fixed base point $x_0\in\matH^3$. It follows from the definitions that $|\alpha(t)|$ equals the measure of $S_0$ for sufficiently big $t$. Moreover the set $S_0$ is in natural correspondence with the set of all isometries that send $x_0$ to some point in $D(x_0)$: its volume does not depend on $t$. 

To estimate $|\alpha(t)^{\rm bad}|$ we fix $g_0\in \Gamma$ so that $D(g_0x_0)\subset U_0$. Let $S^{\rm bad}\subset S(t)$ be the set of all bad $t$-simplices with first vertex in $D(g_0x_0)$. If $t$ is sufficiently big, the volume of $S^{\rm bad}$ is bigger than a constant independent of $t$ (exercise). 

We have proved that $|\alpha(t)^{\rm bad}| / |\alpha(t)|>C>0$ independently of $t$. We may suppose that $\alpha(t)$ represents $[M]$ up to renormalising. The map $f\colon M \to N$ has degree $\pm 1$ and it sends $\alpha(t)$ to a class 
$$f_*(\alpha(t)) = \sum_{\sigma\in \Sigma} \lambda_\sigma(t) \cdot (f\circ\sigma)^{\rm st} $$
representing $\pm [N]$. Since a $C$-portion of $\alpha(t)$ is bad, a $C$-portion of simplices in $f_*(\alpha(t))$ has volume smaller than $v_3-\delta$ and hence
$$\Vol(N) = \left|\int_{f_*(\alpha(t))} \omega\right| < |\alpha(t)| ( (1-C)v_3 + C(v_3 - \delta)) = |\alpha(t)| (v_3 -\delta C).$$
Since this holds for all $t$ and $|\alpha(t)| \to \|M\|$ we get
$$\Vol(N) < \|M\| (v_3 - \delta C) = \Vol (M) - \delta C \cdot \|M\|.$$
Corollary \ref{stesso:volume:cor} gives $\Vol(M) = \Vol(N)$, a contradiction.
\end{proof}

\begin{prop} \label{riflessione:prop}
Every ideal triangle in $\matH^3$ is the face of precisely two regular ideal tetrahedra.
\end{prop}
\begin{proof}
Pick the line $l$ orthogonal to the barycenter of the triangle: the vertex of a regular ideal tetrahedron must be an endpoint of $l$. To prove that these vertices give regular ideal tetrahedra, note that all ideal triangles in $\matH^3$ are isometric, so one concrete example suffices. 
\end{proof}

We turn back to Mostow rigidity.

\begin{prop}
Let $f\colon M \to N$ be a smooth homotopic equivalence between closed hyperbolic orientable 3-manifolds.  
The restriction $\tilde f|_{\partial \matH^3}\colon \partial \matH^3 \to \partial \matH^3$ is the trace of an isometry $\psi\colon \matH^3 \to \matH^3$.
\end{prop}
\begin{proof}
Let $v_0,\ldots, v_3\in\partial \matH^n$ be vertices of a regular ideal tetrahedron $\Delta$. By Lemma \ref{regolare:lemma} the lift $\tilde f$ sends them to the vertices of some regular ideal tetrahedron, and let $\psi$ be the unique isometry of $\matH^3$ such that $\psi (v_i) = \tilde f(v_i)$ for all $i$. 

By Proposition \ref{riflessione:prop} there is a unique point $v_4\neq v_3$ such that $v_0, v_1, v_2, v_4$ are the vertices of an ideal regular tetrahedron, and $\psi (v_4)$ is the unique point other than $\psi(v_3)$ such that $v_0,v_1,v_2, \psi(v_4)$ are the vertices of an ideal regular tetrahedron. By Lemma \ref{regolare:lemma} we also have $\tilde f(v_4) = \psi(v_4)$.

If we mirror $\Delta$ along its faces iteratively we get a tessellation of $\matH^3$ via regular ideal tetrahedra, whose ideal vertices form a dense subset of $\partial \matH^3$. By iterating this argument in all directions the functions $\psi$ and $\tilde f$ coincide on this dense subset and hence on the whole of $\partial \matH^3$.
\end{proof}

We can finally prove Mostow's rigidity theorem.
\begin{teo}[Mostow Rigidity] Let $f\colon M \to N$ be a homotopic equivalence between closed orientable hyperbolic 3-manifolds. The map $f$ is homotopically equivalent to an isometry.
\end{teo}
\begin{proof}
Set $M = \matH^3/_\Gamma$ and $N = \matH^3/_{\Gamma'}$, and pick a lift $\tilde f$. We have
\begin{equation} \label{equivariance:eqn}
\tilde f\circ g = f_*(g)\circ \tilde f \quad \forall g \in \Gamma
\end{equation}
for an isomorphism $f_*\colon \Gamma \to \Gamma'$. We may suppose $f$ smooth. The boundary extension of $\tilde f$ is the trace of an isometry $\psi\colon \matH^3 \to \matH^3$ and hence
\begin{equation} \label{equivariance2:eqn}
\psi\circ g = f_*(g)\circ \psi \quad \forall g \in \Gamma
\end{equation}
holds at $\partial \matH^3$. Both terms in (\ref{equivariance2:eqn}) are isometries, and isometries are determined by their boundary traces: hence (\ref{equivariance2:eqn}) holds for all points in $\matH^3$. Therefore $\psi$ descends to an isometry 
$$\psi\colon M \to N.$$
A homotopy between $f$ and $\psi$ may be constructed from a convex combination of $\tilde f$ and $\psi$ in $\matH^n$, which is also $\Gamma$-equivariant and hence descends.
\end{proof}

\subsection{Consequences of Mostow rigidity}
The most important consequence is that the entire geometry of a closed hyperbolic $3$-manifold is a topological invariant:  numerical quantities like the volume of the manifold, its geodesic spectrum, etc.~depend only on the topology of the manifold.
We single out another application. 

\begin{cor} Let $M$ be a closed orientable hyperbolic 3-manifold. The natural map
$$\Iso (M) \to \Out(\pi_1(M))$$
is an isomorphism.
\end{cor}
\begin{proof}
We already know that it is injective by Proposition \ref{injective:prop}. We prove that it is surjective:
every automorphism of $\pi_1(M)$ is represented by a homotopy equivalence since $M$ is aspherical (see Corollary \ref{aspherical:cor}), which is in turn homotopic to an isometry by Mostow's rigidity.
\end{proof}
We note that this is false in dimension $n=2$, where $\Iso(S)$ is finite and $\Out(\pi_1(S))$ is infinite. 

\subsection{Orbifolds and cone manifolds}
While by Mostow's theorem an orientable closed 3-manifold $M$ can have at most one hyperbolic structure, it may have plenty of hyperbolic cone manifold or orbifold structures, typically distinguished by their singular sets and their cone angles. For instance $M=S^3$ has plenty of such structures, as we will see.

For the moment we note that the volumes of these hyperbolic orbifold structures may be arbitrarily big, but not arbitrarily small:

\begin{prop} \label{Vol:M:O:prop}
Let a closed orientable 3-manifold $M$ be the underlying space of a hyperbolic orbifold $O$. We have
$$\Vol(O) \geqslant v_3\|M\|.$$
\end{prop}
\begin{proof}
By Selberg's Lemma there is a degree-$d$ orbifold covering $N \to O$ which is a closed hyperbolic manifold. We have $\Vol(N) = d\Vol(O)$, and since the covering is a degree-$d$ map $N \to M$ we can apply Corollary \ref{degree:d:cor} and get $\|M\| \leqslant \frac{\|N\|}d = \frac{\Vol(N)}{dv_3} = \frac{\Vol(O)}{v_3}.$
\end{proof}

\subsection{References}
Most of the proofs presented in this chapter were taken from Benedetti -- Petronio \cite{BP}, and are also contained in Thurson's notes \cite{Th}. The proof of Mostow's rigidity presented here was proposed by Gromov and applies to every dimension $n\geqslant 3$, because ideal regular simplexes have indeed the maximum volume by Haagerup and Munkholm \cite{HM}. 
